%package list
\documentclass{article}
\usepackage{amsmath}
\usepackage[top=3cm, bottom=3cm, outer=3cm, inner=3cm]{geometry}
\usepackage{multicol}
\usepackage{graphicx}
\usepackage{url}
%\usepackage{cite}
\usepackage{hyperref}
\usepackage{array}
%\usepackage{multicol}
\newcolumntype{x}[1]{>{\centering\arraybackslash\hspace{0pt}}p{#1}}
\usepackage{natbib}
\usepackage{pdfpages}
\usepackage{multirow}
\usepackage[normalem]{ulem}
\useunder{\uline}{\ul}{}
\usepackage{svg}
\usepackage{xcolor}
\usepackage{minted}

%\usepackage{booktabs}
\usepackage{caption}
\usepackage{subcaption}
\usepackage{float}
\usepackage{array}

\newcolumntype{M}[1]{>{\centering\arraybackslash}m{#1}}
\newcolumntype{N}{@{}m{0pt}@{}}

%%%%%%%%%%%%%%%%%%%%%%%%%%%%%%%%%%%%%%%%%%%%%%%%%%%%%%%%%%%%%%%%%%%%%%%%%%%%
%%%%%%%%%%%%%%%%%%%%%%%%%%%%%%%%%%%%%%%%%%%%%%%%%%%%%%%%%%%%%%%%%%%%%%%%%%%%
\newcommand{\itemEmail}{ccapatinta@unsa.edu.pe}
\newcommand{\itemStudent}{Capatinta Almanza Camilo B. A.}
\newcommand{\itemCourse}{Lab. Análisis y Diseño de Algoritmos}
\newcommand{\itemCourseCode}{1702231}
\newcommand{\itemSemester}{IV}
\newcommand{\itemUniversity}{Universidad Nacional de San Agustín de Arequipa}
\newcommand{\itemFaculty}{Facultad de Ingeniería de Producción y Servicios}
\newcommand{\itemDepartment}{Departamento Académico de Ingeniería de Sistemas e Informática}
\newcommand{\itemSchool}{Escuela Profesional de Ingeniería de Sistemas}
\newcommand{\itemAcademic}{2026 - B}
\newcommand{\itemInput}{12/11/2026}
\newcommand{\itemOutput}{16/11/2026}
\newcommand{\itemPracticeNumber}{10}
\newcommand{\itemTheme}{Métodos Evolutivos y Metaheurísticos.}
%%%%%%%%%%%%%%%%%%%%%%%%%%%%%%%%%%%%%%%%%%%%%%%%%%%%%%%%%%%%%%%%%%%%%%%%%%%%
%%%%%%%%%%%%%%%%%%%%%%%%%%%%%%%%%%%%%%%%%%%%%%%%%%%%%%%%%%%%%%%%%%%%%%%%%%%%

\usepackage[english,spanish]{babel}
\usepackage[utf8]{inputenc}
\AtBeginDocument{\selectlanguage{spanish}}
\renewcommand{\figurename}{Figura}
\renewcommand{\refname}{Referencias}
\renewcommand{\tablename}{Tabla} %esto no funciona cuando se usa babel
\AtBeginDocument{%
	\renewcommand\tablename{Tabla}
}

\usepackage{fancyhdr}
\pagestyle{fancy}
\fancyhf{}
\setlength{\headheight}{40.51407pt}
\renewcommand{\headrulewidth}{1pt}
\renewcommand{\footrulewidth}{1pt}
\fancyhead[L]{\raisebox{-0.2\height}{\includegraphics[width=3cm]{../img/logo_episunsa.png}}}
\fancyhead[C]{\fontsize{7}{7}\selectfont	\itemUniversity \\ \itemFaculty \\ \itemDepartment \\ \itemSchool \\ \textbf{\itemCourse}}
\fancyhead[R]{\raisebox{-0.2\height}{\includegraphics[width=1.2cm]{../img/logo_abet}}}
\fancyfoot[L]{Estudiante Capatinta C.}
\fancyfoot[C]{\itemCourse}
\fancyfoot[R]{Página \thepage}

% para el codigo fuente

\usepackage{color, colortbl}
\definecolor{dkgreen}{rgb}{0,0.6,0}
\definecolor{gray}{rgb}{0.5,0.5,0.5}
\definecolor{mauve}{rgb}{0.58,0,0.82}
\definecolor{codebackground}{rgb}{0.95, 0.95, 0.92}
\setminted{%
  fontsize=\small,
  breaklines,
  linenos,
  numbersep=6pt,
  frame=lines,
  framesep=3mm,
  baselinestretch=1,
  bgcolor=codebackground,
  tabsize=3
}
\definecolor{tablebackground}{rgb}{0.8, 0, 0}

\begin{document}
	
	\vspace*{10px}
	
	\begin{center}
		\parbox{0.8\textwidth}{\centering\fontsize{17}{20}\selectfont\textbf{Informe de laboratorio \itemPracticeNumber}}
	\end{center}
	\begin{center}
		\parbox{0.8\textwidth}{\centering\textbf{\Large Tema: \itemTheme}}
	\end{center}
	%\vspace*{0.5cm}	

	\begin{flushright}
		\begin{tabular}{|M{2.5cm}|N|}
			\hline 
			\rowcolor{tablebackground}
			\color{white} \textbf{Nota}  \\
			\hline 
			     \\[30pt]
			\hline 			
		\end{tabular}
	\end{flushright}	

	\begin{table}[h!]
		\begin{tabular}{|x{4.7cm}|x{4.8cm}|x{4.8cm}|}        
			\hline 
			\rowcolor{tablebackground}
			\color{white} \textbf{Estudiante} & \color{white}\textbf{Escuela}  & \color{white}\textbf{Asignatura}   \\
			\hline 
			{\itemStudent \par \itemEmail} & \itemSchool & {\itemCourse \par Semestre: \itemSemester \par Código: \itemCourseCode}     \\
			\hline 			
		\end{tabular}
	\end{table}		
	
	\begin{table}[H]
		\begin{tabular}{|x{4.7cm}|x{4.8cm}|x{4.8cm}|}
			\hline 
			\rowcolor{tablebackground}
			\color{white}\textbf{Laboratorio} & \color{white}\textbf{Tema}  & \color{white}\textbf{Duración}   \\
			\hline 
			\itemPracticeNumber & \itemTheme & 02 horas   \\
			\hline 
		\end{tabular}
	\end{table}
	
	\begin{table}[H]
		\begin{tabular}{|x{4.7cm}|x{4.8cm}|x{4.8cm}|}
			\hline 
			\rowcolor{tablebackground}
			\color{white}\textbf{Semestre académico} & \color{white}\textbf{Fecha de inicio}  & \color{white}\textbf{Fecha de entrega}   \\
			\hline 
			\itemAcademic & \itemInput &  \itemOutput  \\
			\hline 
		\end{tabular}
	\end{table}
	
\section{Ejercicios propuestos por el docente}

\subsection{Búsqueda Local Iterada (ILS)}

Problema: dado un TSP de 12 ciudades con coordenadas {0,0}, {3,14}, {18,5}, {8,7}, {6,2}, {10,0}, {4,4},
{1,13}, {19,11}, {8,14}, {2,17}, {7,18}, encontrar la ruta más corta que visita todas las ciudades y regresa a la
ciudad 0. Con 12 ciudades la fuerza bruta ya sería lenta (11! rutas), así que usamos metaheurísticas. El
óptimo, verificado con un algoritmo exacto, es 72.21; úsalo para medir qué tan buenas son las soluciones.

\inputminted{cpp}{../e-p/busqueda_local_iterada.cpp}

\begin{figure}[H]
    \centering
    \includegraphics[width=0.95\textwidth]{../img/ILS.png}
    \caption{Ejecución de \texttt{busqueda\_local\_iterada.cpp}.}
    \label{fig:ILS}
\end{figure}

\begin{itemize}
	\item Explicación1
\end{itemize}

\begin{figure}[H]
    \centering
    \includegraphics[width=0.95\textwidth]{../img/busqueda_local_iterada.png}
    \caption{Comportamiento observada de búsqueda local iterada.}
    \label{fig:ILS_diagrama}
\end{figure}

\begin{itemize}
	\item Explicación1
\end{itemize}


\subsection{Algoritmo Iterativo Goloso (IG)}
\inputminted{cpp}{../e-p/algoritmo_iterativo_goloso.cpp}

\begin{figure}[H]
    \centering
    \includegraphics[width=0.95\textwidth]{../img/IG.png}
    \caption{Ejecución de \texttt{algoritmo\_iterativo\_goloso.cpp}.}
    \label{fig:IG_diagrama}
\end{figure}

\begin{itemize}
	\item Explicación1
\end{itemize}

\begin{figure}[H]
    \centering
    \includegraphics[width=0.9\textwidth]{../img/algoritmo_iterativo_goloso.png}
    \caption{Comportamiento observado del algoritmo iterativo goloso.}
    \label{fig:IG}
\end{figure}

\begin{itemize}
	\item Explicación1
\end{itemize}

\section{Ejercicios/Problemas propuestos}

\subsection{Header compartido: \texttt{tsp\_meta.h}}
\inputminted{cpp}{../e-n/tsp.h}

El archivo \texttt{tsp\_meta.h} define las estructuras y funciones comunes utilizadas
para resolver el Problema del Agente Viajero (TSP) mediante ILS e IG.

La estructura \texttt{Punto} representa las coordenadas de una ciudad, mientras que
\texttt{Matriz} almacena las distancias entre todas las ciudades y \texttt{Ruta}
contiene los índices que indican el orden de visita.

La función \texttt{instancia12} proporciona las 12 ciudades del laboratorio.
\texttt{matrizDist} calcula las distancias euclidianas mediante \texttt{hypot},
mientras que \texttt{costo} obtiene el costo total del recorrido cerrado.
\texttt{generarPuntos} permite crear instancias aleatorias reproducibles mediante
una semilla.

\texttt{vecinoMasCercano} construye una solución inicial desde la ciudad 0,
seleccionando en cada paso la ciudad no visitada más cercana. Posteriormente,
\texttt{dosOpt} mejora esta ruta invirtiendo tramos siempre que la variación
$\Delta$ del costo sea negativa. El proceso se repite hasta alcanzar un óptimo local,
es decir, hasta que ninguna inversión 2-opt considerada produzca una mejora.

En ILS, \texttt{doubleBridge} perturba la ruta dividiéndola en tres puntos y
reordenando los segmentos $A\,B\,C\,D$ como $A\,C\,B\,D$. Después se vuelve a
aplicar \texttt{dosOpt}. La función \texttt{ils} repite este proceso con un criterio
de aceptación estricta: conserva una solución nueva solo si mejora la solución actual.
Además, almacena por separado la mejor ruta obtenida durante toda la ejecución.

En IG, \texttt{ig} selecciona aleatoriamente \texttt{nDestruir} ciudades y las elimina
de la ruta actual. Luego, las reinserta una a una en la posición que produce el menor
incremento de costo. Este mecanismo permite conservar parte de la estructura buena,
destruir una sección y reconstruir una nueva solución. Como en ILS, conserva únicamente
una ruta nueva si mejora la solución actual y puede aplicar búsqueda local mediante
\texttt{dosOpt}.

Finalmente, \texttt{optimoHeldKarp} determina el costo óptimo mediante programación
dinámica sobre subconjuntos de ciudades. Al tener complejidad exponencial, solo es
práctico para instancias pequeñas y se utiliza como referencia para evaluar la calidad
de las soluciones obtenidas por las metaheurísticas.

En conjunto, \texttt{tsp\_meta.h} centraliza la instancia, el cálculo de distancias,
la construcción de rutas, la búsqueda local, las perturbaciones, las metaheurísticas
ILS e IG y el cálculo del óptimo exacto.

\subsection{Ejercicio 1}
Ejecuta el ILS del ejercicio resuelto con 10 semillas distintas (por ejemplo 1 a 10) y registra para cada una el costo final y la iteración en que se encontró. ¿En cuántas semillas se llega al óptimo 72.21? ¿Qué te dice eso sobre la confiabilidad de un algoritmo aleatorio?
\inputminted{cpp}{../e-n/e1.cpp}

\begin{figure}[H]
    \centering
    \includegraphics[width=0.8\textwidth]{../img/e1.png}
    \caption{Ejecución del ejercicio 1}
    \label{fig:e1}
\end{figure}

\subsection{Ejercicio 2}
Modifica el IG para probar d = 1, 2, 3, 5 y 8 ciudades destruidas. Para cada valor, ejecútalo con 10 semillas y reporta el costo promedio. ¿Qué pasa con d muy pequeño? ¿Y con d muy grande? Relaciona tus resultados con los conceptos de intensificación y diversificación.
\inputminted{cpp}{../e-n/e2.cpp}

\begin{figure}[H]
    \centering
    \includegraphics[width=0.8\textwidth]{../img/e2.png}
    \caption{Ejecución del ejercicio 2}
    \label{fig:e2}
\end{figure}

\subsection{Ejercicio 3}
Cambia el criterio de aceptación del ILS: además de aceptar soluciones mejores, acepta una solución peor con probabilidad 5\%. Compara el resultado con el criterio ``solo si mejora'' usando varias semillas. ¿Ayuda aceptar soluciones peores a escapar de los valles? Guarda siempre la mejor solución encontrada, aunque la actual empeore.
\inputminted{cpp}{../e-n/e3.cpp}

\begin{figure}[H]
    \centering
    \includegraphics[width=0.8\textwidth]{../img/e3.png}
    \caption{Ejecución del ejercicio 3}
    \label{fig:e3}
\end{figure}

\subsection{Ejercicio 4}
Genera un TSP aleatorio de 50 ciudades (con semilla fija) y compara cuatro métodos: Nearest Neighbor, NN+2-opt, ILS e IG con 500 iteraciones cada uno. Reporta costo y tiempo de ejecución en una tabla. ¿Cuánto mejora cada metaheurística respecto a NN+2-opt y a qué costo de tiempo?
\inputminted{cpp}{../e-n/e4.cpp}

\begin{figure}[H]
    \centering
    \includegraphics[width=0.8\textwidth]{../img/e4.png}
    \caption{Ejecución del ejercicio 4}
    \label{fig:e4}
\end{figure}

\subsection{Ejercicio 5}
Aplica la metaheurística que mejor te funcionó a un caso real: el recorrido de entrega por 15 a 20 puntos de tu ciudad (o capitales del Perú) con coordenadas reales o aproximadas. Muestra la ruta final y compárala con la ruta que obtendrías con solo Nearest Neighbor. Justifica por qué elegiste ILS o IG.
\inputminted{cpp}{../e-n/e5.cpp}

\begin{figure}[H]
    \centering
    \includegraphics[width=0.8\textwidth]{../img/e5.png}
    \caption{Ejecución del ejercicio 5}
    \label{fig:e5}
\end{figure}

\section{Cuestionario}

\begin{itemize}
    \item ¿Qué es un óptimo local? ¿Por qué 2-opt, por sí solo, no puede salir de él aunque exista una ruta mejor?
    \item ¿Qué diferencia hay entre intensificación y diversificación? Indica qué componente de ILS y qué componente de IG se encargan de cada una.
    \item En ILS y en IG la solución nueva solo se acepta si mejora la actual. ¿Qué ventaja y qué riesgo tiene este criterio de aceptación?
    \item ¿Qué pasaría en ILS si la perturbación fuera demasiado débil? ¿Y si fuera tan fuerte que cambiara casi toda la ruta? Explica con el paisaje de soluciones de la Figura 1.
    \item Según la Figura 7, ILS baja más rápido por iteración que IG. ¿Significa eso que ILS es siempre mejor? Considera el costo de cada iteración en ambos algoritmos.
\end{itemize}

\section{Conclusiones}

\end{document}